\section{Chat Template}
\label{sec:post-chat-template}

The \kimi{3} chat template is redesigned around three goals. The first is \emph{extensibility}: new capabilities should be introduced through backward-compatible message formats rather than template revisions, so that a single template serves the entire model generation. The second is a \emph{low alignment tax}: the format should be learnable with minimal supervised data, supporting a pipeline in which a lightly fine-tuned pre-trained model can proceed directly to reinforcement learning. The third is \emph{decoding friendliness}: the structure should admit simple encoders, streaming parsers, and grammar-constrained enforcers. To these ends, the template adopts XTML (eXtensible Token Markup Language), an XML-like markup in which the angle-bracket syntax is replaced by three reserved special tokens: \texttt{[open]}, \texttt{[sep]} and \texttt{[close]}, with an additional \texttt{[end\_of\_msg]} token as the generation stop marker. An element \texttt{[open]tag attr="value"[sep] \dots\ [close]tag[sep]} is isomorphic to its XML counterpart, but every structural boundary is an explicit special token, which removes tokenization ambiguity at element boundaries and simplifies constrained decoding.

\begin{figure}[t]
    \centering
    \definecolor{ctfigink}{HTML}{1C1C1E}
\definecolor{ctfiggray}{HTML}{777777}
\definecolor{ctfiggrid}{HTML}{D6D6D8}
\definecolor{ctfigred}{HTML}{B92622}
\definecolor{ctfigredfill}{HTML}{FBF1F0}
\definecolor{ctfigblue}{HTML}{1773E6}
\definecolor{ctfigbluefill}{HTML}{EDF4FD}

\begin{adjustbox}{width=\textwidth,center}
        \begin{tikzpicture}[
                        x=1pt,
                        y=1pt,
                        outer sep=0pt,
                        paneltitle/.style={font=\fontsize{9.2}{10.8}\selectfont, text=ctfigink, align=center},
                        zonelabel/.style={font=\fontsize{7.2}{8.4}\selectfont, text=ctfiggray},
                        zone/.style={draw=ctfiggrid, line width=0.5pt, rounded corners=2.5pt, fill=white},
                        chip/.style={draw=ctfigink!75, line width=0.5pt, rounded corners=1.5pt,
                                        fill=white, font=\fontsize{6.9}{8}\selectfont, text=ctfigink,
                                        inner xsep=3pt, inner ysep=2.2pt},
                        optchip/.style={chip, draw=ctfigblue!80, fill=ctfigbluefill},
                        hotchip/.style={chip, draw=ctfigred, fill=ctfigredfill, line width=0.7pt},
                        expchip/.style={chip, dashed, draw=ctfiggray},
                        chan/.style={line width=0.5pt, rounded corners=2pt},
                        chanthink/.style={chan, draw=ctfiggray!75, fill=white},
                        chanresp/.style={chan, draw=ctfigblue!65, fill=ctfigbluefill},
                        chantools/.style={chan, draw=ctfigred!60, fill=ctfigredfill},
                        callcard/.style={draw=ctfigink!55, line width=0.45pt, rounded corners=1.8pt, fill=white},
                        mono/.style={font=\scriptsize\ttfamily, text=ctfigink},
                        textline/.style={draw=ctfiggrid, line width=1.7pt, line cap=round},
                        flow/.style={draw=ctfiggray, line width=0.55pt,
                                                -{Latex[length=3.8pt,width=3pt]}},
                        zoom/.style={draw=ctfiggray!85, line width=0.45pt, dashed},
                ]
                \path[use as bounding box] (0,-4) rectangle (560,205);

                \node[paneltitle, anchor=north] at (82,203)
                {(a) Context layout};

                \draw[zone] (6,150) rectangle (158,184);
                \node[zonelabel, anchor=north] at (82,182) {global option messages};
                \node[optchip] at (47,162) {tool-declare};
                \node[optchip] at (110,162) {thinking-effort};

                \draw[zone] (6,80) rectangle (158,140);
                \node[zonelabel, anchor=north] at (82,138) {input messages};
                \node[chip] (m-sys) at (48,122) {system};
                \node[chip] (m-usr) at (112,122) {user};
                \node[chip] (m-tool) at (48,108) {tool};
                \node[hotchip] (m-asst) at (112,108) {assistant};
                \node[expchip] (m-dyn) at (82,94) {dynamic tool-declare};

                \draw[zone] (6,38) rectangle (158,72);
                \node[zonelabel, anchor=north] at (82,70) {one-shot option messages};
                \node[optchip] at (46,50) {tool-choice};
                \node[optchip] at (112,50) {response-format};

                \draw[flow] (82,147) -- (82,143);
                \draw[flow] (82,77) -- (82,75);

                \node[chip] (genprefix) at (43.5,24) {[open]think[sep]};
                \node[zonelabel] at (79.7,24) {/};
                \node[chip] at (115.8,24) {[open]response[sep]};
                \node[zonelabel, anchor=north] at (82,15) {generation prefix};
                \draw[flow] (82,36) -- (82,30);

                \node[paneltitle, anchor=north] at (270,203)
                {(b) Assistant message};

                \draw[zoom] (m-asst.north east) -- (176,180);
                \draw[zoom] (m-asst.south east) -- (176,22);

                \draw[zone] (180,18) rectangle (360,184);
                \node[mono, anchor=north west] at (186,181) {[open]message role="assistant"[sep]};

                \draw[chanthink] (188,124) rectangle (352,170);
                \node[mono, anchor=north west] at (194,167) {[open]think[sep]};
                \draw[textline] (198,151) -- (272,151);
                \draw[textline] (198,143.5) -- (256,143.5);
                \node[mono, text=ctfiggray, anchor=south west] at (194,127) {[close]think[sep]};

                \draw[chanresp] (188,80) rectangle (352,118);
                \node[mono, anchor=north west] at (194,115) {[open]response[sep]};
                \draw[textline] (198,101) -- (272,101);
                \draw[textline] (198,93.5) -- (256,93.5);
                \node[mono, text=ctfiggray, anchor=south west] at (194,83) {[close]response[sep]};

                \draw[chantools] (188,32) rectangle (352,74);
                \node[mono, anchor=north west] at (194,71) {[open]tools[sep]};
                \node[mono, text=ctfiggray] at (270,53) {$\cdots$};
                \node[mono, text=ctfiggray, anchor=south west] at (194,35) {[close]tools[sep]};

                \node[mono, text=ctfiggray, anchor=south west] (msgclose) at (186,21.5) {[close]message[sep]};
                \node[chip, right=3.5pt of msgclose] {[end\_of\_msg]};

                \node[paneltitle, anchor=north] at (465,203)
                {(c) Tools channel};

                \draw[zoom] (352,72) -- (371,180);
                \draw[zoom] (352,34) -- (371,40);

                \draw[chantools] (373,34) rectangle (559,184);
                \node[mono, anchor=north west] at (379,181) {[open]tools[sep]};

                \draw[callcard] (379,118) rectangle (555,170);
                \node[mono, align=left, anchor=north west, inner xsep=2.5pt, inner ysep=2pt]
                at (384,167.5) {[open]call tool="python" index="1"[sep]\\[0.5pt]
                        {\color{ctfiggray}\fontsize{6.0}{7.2}\selectfont [open]argument key="code" type="string"[sep]}\\[0.5pt]
                        {\color{ctfiggray}\fontsize{6.0}{7.2}\selectfont \hspace{1.2em}$\cdots$}\\[0.5pt]
                        {\color{ctfiggray}\fontsize{6.0}{7.2}\selectfont [close]argument[sep]}\\[0.5pt]
                        {\color{ctfiggray}[close]call[sep]}};

                \draw[callcard] (379,56) rectangle (555,110);
                \node[mono, align=left, anchor=north west, inner xsep=2.5pt, inner ysep=2pt]
                at (384,107.5) {[open]call tool="search" index="2"[sep]\\[0.5pt]
                        {\color{ctfiggray}\fontsize{6.0}{7.2}\selectfont [open]argument key="options" type="object"[sep]}\\[0.5pt]
                        {\color{ctfiggray}\fontsize{6.0}{7.2}\selectfont \hspace{1.2em}\{"timeout": 150\}}\\[0.5pt]
                        {\color{ctfiggray}\fontsize{6.0}{7.2}\selectfont [close]argument[sep]}\\[0.5pt]
                        {\color{ctfiggray}[close]call[sep]}};

                \node[mono, text=ctfiggray, anchor=north west] at (379,48) {[close]tools[sep]};
        \end{tikzpicture}
\end{adjustbox}

    \caption{Structure of the \kimi{3} chat template.
        \textbf{(a)} Context layout: global option messages precede the input messages, while one-shot option messages follow them, so that per-request options leave the history KV cache intact; dynamically loaded tools are injected mid-session as input option messages (dashed).
        \textbf{(b)} Anatomy of an assistant message: the body is organized into \texttt{think}, \texttt{response}, and \texttt{tools} channels.
        \textbf{(c)} Expansion of the \texttt{tools} channel: parallel tool calls are indexed so that tool results can be matched to their calls, and arguments are typed.}
    \label{fig:chat-template}
\end{figure}

\paragraph{Messages and zones}
The top-level unit of the context is the message, and messages fall into two categories by origin (Fig.~\ref{fig:chat-template}a). \emph{Input messages} serialize the \texttt{messages} field of the request, covering the familiar system, user, assistant, and tool roles. \emph{Option messages} translate request options into instructions that the model reads in context, and their placement reflects their scope. \emph{Global options}---the tool declaration (\texttt{type="tool-declare"}) and the reasoning-effort setting---appear before all input messages: they govern the whole session and rarely change, so modifying them invalidates the KV cache anyway. \emph{One-shot options} (\texttt{tool\_choice}, \texttt{response\_format}) are appended after the input messages, so that per-request changes leave the history KV cache intact. A third kind, the \emph{input option message}, is interleaved with input messages to supplement or override a global option mid-session. This mechanism supports \emph{dynamically loaded tools}: tools retrieved or loaded during a conversation are announced through an additional tool-declare message, after which the model's available toolset expands without rebuilding the preceding context.

\paragraph{Channels}
The body of an assistant message is organized into \emph{channels}, a concept inspired by OpenAI's Harmony response format~\citep{openai2025harmony}: \texttt{think} carries the reasoning trace, \texttt{response} the user-visible answer, and \texttt{tools} the tool calls (Fig.~\ref{fig:chat-template}b). The two generation modes are selected purely through the generation prefix---\texttt{[open]think[sep]} for thinking mode and \texttt{[open]response[sep]} for instruct mode---rather than through separate templates. \kimi{3} supports only \emph{preserved thinking}: in thinking mode, the think channel is always retained in the history---kept even when its content is empty---so that the model observes a consistent message structure across turns; in instruct mode, historical messages contain only the response and tools channels.

\paragraph{Tool calling}
Within the tools channel, each call carries \texttt{tool} and \texttt{index} attributes; the index numbers parallel calls within a message, and each tool-result message repeats the same \texttt{tool}/\texttt{index} pair and follows the order of its call, so that results are unambiguously associated with calls. Arguments are typed: string arguments appear as raw text, while values of other JSON types are compactly serialized. Free-form text such as code is therefore a first-class citizen rather than an escaped JSON string. A pure-JSON fallback block covers inputs whose arguments cannot be decomposed into typed argument blocks; it occurs only in input tokens, never in model outputs, and its loss is masked during training.

\paragraph{Reasoning effort and options}
Reasoning effort is exposed as a global option message of type \texttt{thinking-effort}, inserted after the tool declaration and before the input messages. Instead of modifying the generation prefix or exposing a token budget, the message states the requested level in natural language and acts as a generation-constraint instruction. The schema reserves four levels (\texttt{low}, \texttt{medium}, \texttt{high}, and \texttt{max}), of which \kimi{3} supports a subset. This representation decouples the effort interface from the template syntax, and it aligns directly with the effort-conditioned training described in \S\ref{sec:post-sft} and \S\ref{sec:post-rl}.

More broadly, this is the common implementation of all option messages: \texttt{tool\_choice}, \texttt{response\_format}, and \texttt{thinking-effort} are each translated into a short natural-language instruction placed in context, rather than into dedicated special syntax. Because the pre-trained model already follows such instructions well, new options can be introduced with little or no additional training---a direct embodiment of the low-alignment-tax design principle stated above.
